The textbook prediction that an import tariff appreciates the tariffing country's currency comes from a two-country world, in which the bilateral exchange rate, the effective exchange rate and the terms of trade are the same object. With many countries they come apart. We characterize the long-run exchange-rate response to arbitrary tariff vectors in an $N$-country trade model. To first order, the response is the inverse of an exchange-rate adjustment matrix applied to the tariff's trade-balance disturbance at unchanged exchange rates, and a multilateral Marshall--Lerner condition makes that adjustment well behaved. Under partner symmetry, any tariff vector decomposes into \emph{average protection}, which alone moves the nominal effective exchange rate, and \emph{relative treatment}, which alone moves bilateral rates around it.

In a nested-CES model these responses are governed by two margins: diversion toward the domestic economy, $\rho_D$, and substitution across foreign suppliers, $\rho$. A unilateral tariff always appreciates the tariffing country's effective exchange rate, yet it depreciates the currency against every untariffed bystander when $\rho>N\rho_D$, which requires that the two elasticities differ. Under retaliation the threshold falls to $\rho_D$, and both belligerents depreciate in effective terms. The symmetric $N$-country economy reduces exactly to a three-country economy with a composite bystander. Under partner asymmetry, a uniform tariff still appreciates the currency against every partner and raises welfare, but a tariff on a small, low-weight partner can worsen the tariffing country's aggregate terms of trade, and hence its welfare, when foreign suppliers are close substitutes. Read stylistically, the 2025 retaliation phase is consistent with this logic, with a roughly stable dollar against the renminbi and weakness against bystanders, whereas the dollar's effective weakness under unilateral and near-uniform tariffs is not explained by the trade channel.
